\subsection{Préliminaires III : variance, complexes et limites}
\label{note:base3-complexes}
Un foncteur contravariant vérifie $F(vu)=F(u)F(v)$ (8.1.3).
La famille booléenne de parties constructibles en 9.1.2 contient
$X$. L'ouvert non vide contenu dans la partie localement fermée
de 9.2.3 est un ouvert de $X$, pas seulement un ouvert relatif
de cette partie. Le morphisme du complété en 10.2.6 est
$\widehat v$ ; en 10.2.8 le morphisme local induit l'application
des corps résiduels $k\to k_i$.

Dans 11.1.6, le terme $E_r^{n-q(n),q(n)}$ appartient au gradué
d'indice $n-q(n)$ de $E^n$. L'inégalité
$p-r+1+j_0<n=p+q$ de 11.3.3 équivaut à $r>j_0-q+1$.
En 11.4.1, les termes du complexe de cochaînes portent les
indices supérieurs $L^i$ ; les limites de 11.4.6 sont prises
dans la catégorie but $\mathcal C''$.
Les flèches de 11.5.2 ont le terme $f'^{i+1}$ correspondant
au premier complexe. Les troncatures de 11.5.4 approchent le
complexe lorsque $r$ tend vers moins l'infini : c'est
$-r$ qui est assez grand pour un degré total fixé.

Les cobords de degré $n$ en 13.2.3 sont un sous-groupe de
$K^n$. Les foncteurs dérivés de 13.2.4 concernent la limite
projective. Enfin, l'opérateur de 13.6.6 a bidegré $(j,-j)$,
qui conserve le degré total et indique exactement le déplacement
dans la filtration.

\subsection{Préliminaires III, 11 et III, 1.1 : signes et degrés homologiques}
\label{note:base3-signs}
La mise en ordre des sommets d'un simplexe en 11.8.1 porte le
signe $\varepsilon_\pi$ de la permutation. C'est la relation
d'alternance qui rend cette opération compatible avec les bords.
Dans 11.8.2, le signe $(-1)^{h+1}$ dépend du degré $h$ du
simplexe ; le nombre de sommets distincts peut être moindre
lorsque des sommets se répètent. En 11.8.6, les groupes du
complexe de chaînes sont $H_q$. Les modules de 11.9.3 sont
des modules à droite, conformément à la catégorie utilisée
dans la résolution.

L'identité de III, 1.1.7,
$i_{\mathbf f}e_{\mathbf g}+e_{\mathbf g}i_{\mathbf f}
=\langle\mathbf g,\mathbf f\rangle1$, est une identité
d'endomorphismes de $\bigwedge(A^r)$. Les opérateurs contraction
et multiplication extérieure parcourent ses différents degrés ;
l'identité ne se limite pas à son terme $A^r$ de degré un.

\subsection{Préliminaires III, 12 : Ext local et changement de base}
\label{note:base3-ext}
Une restriction des sections va de l'ouvert $V$ vers l'ouvert
$V'\subset V$ (12.2.2). Dans 12.3.3, le voisinage est celui du
point $x$. Pour un module localement libre $\mathcal L$, c'est
le faisceau $\mathcal E\!xt^i(\mathcal L,\mathcal G)$ qui
s'annule pour $i>0$ : localement le foncteur Hom d'un module
libre est exact. Les groupes Ext globaux peuvent rester non
nuls, puisqu'ils comportent la cohomologie de
$\mathcal H\!om(\mathcal L,\mathcal G)$.

Le changement de base plat de 12.3.4 fournit le morphisme
$f^*\mathcal E\!xt^\bullet(\mathcal F,\mathcal G)
\to\mathcal E\!xt^\bullet(f^*\mathcal F,f^*\mathcal G)$.
L'énoncé général n'affirme pas que ce soit un isomorphisme.
Dès le degré zéro, commuter Hom à une extension plate de
scalaires demande des hypothèses de finitude sur son premier
argument ; celles-ci ne figurent pas dans cette assertion générale.

\subsection{III, 1--3 : sections, supports et twists}
\label{note:base3-sections}
Les $g_i$ de 1.2.4 sont des fonctions régulières sur $U$,
donc des éléments de $\Gamma(U,\mathcal O_X)$ ; l'application
de degré zéro est $\Gamma(X,\mathcal G)\to H^0(\mathfrak U,\mathcal G)$.
Dans 1.4.15, la quasi-compacité donne un recouvrement affine
\emph{fini}, et la séparation est l'autre hypothèse utilisée.
La finitude du morphisme structurel n'est pas requise pour ce choix.
Le morphisme examiné dans 1.4.17 est $d_1$.
Si $s$ prolonge déjà $f^ns'$ dans cette preuve, la section
obtenue est $f^nt+u(s)$ : appliquer une nouvelle fois $f^n$
à $s$ compterait deux fois ce facteur.

Le twist de III, 2.1.16 donne le terme symétrique d'indice
$-n-r-1$ dans $f_*(\boldsymbol\omega(-n))$.
En 2.2.5(ii), l'obstruction à l'annulation est
$H^r(X,\mathcal O_X)\ne0$ ; elle n'exige pas d'identifier
ce groupe à l'anneau de base. Les annulations de 2.4.1
concernent les degrés $p>0$, pas les sections de degré zéro.
Dans le dévissage de 3.1.2, le module test $\mathcal G$ a
support exactement $Y$ et germe générique de dimension un.
Le contrôle du support permet de réduire ensuite la différence
à un fermé propre de $Y$ ; la seule condition au point générique
ne contrôle pas d'éventuels supports supplémentaires ailleurs.

\subsection{III, 2.2.4--2.2.5 : exactitude et support propre}
\label{note:base3-propersupport}
Les lectures de III, 2.2.4 et 2.2.5(iii) sont celles des
errata des auteurs publiés dans III, seconde partie, pp.~221--222,
Errata III, nos~26--27.

Sous les hypothèses de III, 2.2.4, posons
$\mathcal G_1=\operatorname{Im}(\mathcal F\to\mathcal G)
=\ker(\mathcal G\to\mathcal H)$.
La condition pertinente est la propreté du support de
$\mathcal G_1$. L'exactitude à gauche de $f_*$ ramène alors
l'énoncé à la surjectivité éventuelle de
$f_*\mathcal F(n)\to f_*\mathcal G_1(n)$.

Dans la compactification projective $i:X\hookrightarrow Z$
de la démonstration, un faisceau cohérent $\mathcal G$ de
support propre se prolonge en $i_*\mathcal G$, cohérent et
de même support. Un prolongement cohérent $\mathcal F'$ de
$\mathcal F$ existe par I, 9.4.3. La surjection sur $X$
se recolle avec l'application nulle hors du support de
$\mathcal G$ : elle donne $\mathcal F'\to i_*\mathcal G$.
Pour la base affine, l'annulation de Serre sur $Z$ donne
la surjectivité sur les sections pour les twists assez grands.
La restriction des sections de $i_*\mathcal G(n)$ à $X$
est bijective, d'où la surjectivité désirée sur $X$.

Le même argument sur le support propre donne
$R^if_*\mathcal H(n)=0$ pour $i>0$ et $n$ assez grand.
Si $u:\mathcal F\to\mathcal G$ a noyau et conoyau de supports
propres, on applique les suites longues aux deux suites courtes
définies par son image. Pour le conoyau, la surjectivité
$f_*\mathcal G(n)\to f_*\operatorname{Coker}(u)(n)$ découle
en outre de 2.2.4. Cette surjectivité traite aussi le degré un,
et l'on obtient les isomorphismes
$R^if_*\mathcal F(n)\simeq R^if_*\mathcal G(n)$ pour tout $i>0$.

\subsection{III, 3.4.4 : gradué associé et multiplication}
\label{note:base3-associatedgraded}
Les classes de produits donnent des surjections
\[
 (\mathfrak J^k/\mathfrak J^{k+1})\otimes_{A_0}N_0
 \longrightarrow \mathfrak J^kN/\mathfrak J^{k+1}N,
 \qquad \bar a\otimes\bar n\longmapsto\overline{an},
\]
où $N_0=N/\mathfrak JN$. La surjectivité résulte de la définition
de $\mathfrak J^kN$ par les sommes de produits.
Elle ne donne pas une injectivité sans hypothèse supplémentaire.
Pour $A=k[[t]]$, $\mathfrak J=(t)$ et $N=A/(t)$, le terme
de gauche en degré un est $k$, tandis que celui de droite est zéro.

De même, la multiplication donne une surjection de
$\mathcal O_{X_0}\otimes_{A_0}\operatorname{gr}_{\mathfrak J}(A)$
vers $\operatorname{gr}(\mathcal O_{\mathfrak X})$.
Le module gradué effectivement associé à $N$ est
$M'=\bigoplus_k\mathfrak J^kN/\mathfrak J^{k+1}N$.
Il est quotient de $S'\otimes_{B_0}N_0$ ; la finitude de $N_0$
implique donc sa finitude sur $S'$. La conclusion de la preuve
utilise cette surjection, sans identification injustifiée des gradués.

\subsection{III, 3--4 : bases, corps de fonctions et immersion ouverte}
\label{note:base3-localgeometry}
Dans 3.4.1, le morphisme est $f_n:X_n\to S_n$.
La filtration de 3.4.2 porte l'indice $N_{n,k}$ dans l'ordre
défini. L'image directe de 4.1.5 est un module sur
$\mathcal O_{\widehat Y}$, et la flèche plate de 4.3.4 est $h$.
Dans 4.3.7, la fermeture algébrique du corps $R(Y)$ dans
$R(X)$ est considérée comme une extension de $R(Y)$ ; la
radicialité porte donc sur ce corps de base.
En 4.3.12, la finitude sur $\mathcal O_y$ rend les éléments
algébriques sur son corps des fractions $R(Y)$.
Avec l'hypothèse sur la fermeture algébrique et la normalité
de $\mathcal O_y$, l'inclusion obtenue est une égalité.

Conformément à l’erratum des auteurs (III, seconde partie, p.~222, Errata III, no~28), pour 4.4.9, $X,Y$ sont intègres et localement noethériens,
$Y$ est normal, et $f$ est séparé, localement de type fini et
birationnel, avec fibres finies. Une fibre localement de type
fini sur un corps et finie comme espace topologique est discrète
dans cette situation ; le critère local de 4.4.8 s'applique
aux anneaux locaux de même corps des fractions.
Il donne l'isomorphisme local requis, et la séparation permet
l'immersion ouverte de l'énoncé. Si $f$ est aussi fermé, son
image est ouverte et fermée, non vide et dense dans $Y$
intègre, donc égale à $Y$ ; $f$ est alors un isomorphisme.
Dans l'index final, l'argument du foncteur $h_w$ est l'objet
$Y$ ; $w$ indique le morphisme structural.

\subsection{Préliminaires III et III : renvois}
\label{note:base3-refs}
Le résultat utilisé en 0, 11.8.9 est 11.8.5 ; le morphisme
de 0, 12.1.4 porte le numéro 12.1.4.1 ; la formule de
0, 13.5 est 13.5.5.3. Les renvois de III, 1.4 et 2.3--2.5
concernent respectivement 1.4.15, 2.3.1 et 2.5.2.
Dans III, 3.3, la formule est 3.3.2.1 et le théorème invoqué
est 3.3.1 ; ce dernier est également celui utilisé en 3.4.4.
Les formules de III, 4.1 auxquelles renvoient les preuves
sont 4.1.7.8 et 4.1.3.5.

\subsection{III, 3.4.3 : la base formelle}
\label{note:base3-formalbase}
La base affine dans 3.4.3 est le schéma formel $\mathfrak Y$.
Cette notation distingue la base formelle des schémas obtenus
par réduction modulo les puissances de son idéal de définition.
\par
